% !PRINTED 356
\weberpage{356}
%<<CONT
equal to zero for all values of $(x)$ satisfying the equation $f(x)=0$, and consequently it must be divisible by $f(x)$. We can therefore, if we denote by $\Phi(x,a,b)$ a form of degree $22$, set:

\weq{15}{b_x^6\, D(x,a) - a_x^6\, D(x,b) = f(x)\, \Phi(x,a,b).}

Let us imagine for a moment that a new coordinate system is introduced in (15), in which the lines $a_x, b_x$ are two axes, of which we assume that they do not intersect on the curve $f$; then it follows from the comparison of the right and left sides of the identical equation (15) that no term can occur in $\Phi$ that is not multiplied either by $a_x^6$ or by $b_x^6$, and consequently $\Phi(x)$ has the form

\[
\Phi(x,a,b) = a_x^6\, \Phi_2(x) - b_x^6\, \Phi_1(x),
\]

in which $\Phi_1, \Phi_2$ are forms of order $16$. Accordingly the identical equation (15) takes the form

\[
b_x^6\,[D(x,a) + f(x)\,\Phi_1(x)] = a_x^6\,[D(x,b) + f(x)\,\Phi_2(x)],
\]

and it shows that $D(x,a) + f(x)\,\Phi_1(x)$ is divisible by $a_x^6$. There exists, therefore, a form of degree $14$, $\chi(x,a,b)$, such that

\weq{16}{\frac{D(x,a)}{a_x^6} = \chi(x,a,b) - f(x)\,\frac{\Phi_1(x,a,b)}{a_x^6}.}

If we here substitute for the $a$ and $b$ any special, e.g. rational, values $c, d$, and denote $\chi(x,c,d)$, which then depends only on $x$ and on the coefficients of the form $f(x)$, by $\chi(x)$, then it follows:

\[
\frac{D(x,c)}{c_x^6} = \chi(x) - f(x)\,\frac{\Phi_1(x,c,d)}{c_x^6},
\]

and if we subtract this from (16) and use the relation (15) for $b=c$:

\weq{17}{\chi(x,a,b) - \chi(x) = f(x)\,\Psi(x),}

in which $\Psi(x)$ is a function that in any case can contain no other denominator than a product of powers of $a_x$ and $c_x$. But since $f(x)$ is not divisible by $a_x$ and $c_x$, $\Psi(x)$ must be an integral function, and (16) yields, if we denote by $\Theta(x)$ a new form of $x$ (of degree 16):

\weq{18}{\frac{D(x,a) - f(x)\,\Theta(x)}{a_x^6} = \chi(x).}
